\chapter{相关工作}\label{app:auxiliary related works}

在第 \ref{chap1-sec-related-works} 节中，我们介绍了特征学习理论的研究现状。

In contrast to many machine learning topics, anomaly detection (AD) adopts the \qm{open-world} assumption in which the training and test data are not drawn from the same distribution. 
According to how the test distribution differs from the training distribution, AD can be further categorized into semantic AD \citep{DBLP:conf/aaai/AhmedC20,DBLP:conf/icml/DeeckeRVB21} and non-semantic AD \citep{DBLP:journals/ijcv/BergmannBFSS21,DBLP:journals/corr/abs-2305-15956}.
% \cref{fig-different-types-of-anomalies} provides a toy example of the semantic and non-semantic anomalies.
In this work, we characterize the normal data and both types of anomalies using a unified probabilistic model.

The \qm{open-world} setting of AD is also adopted in some closely related topics, e.g., outlier detection (OD) \citep{DBLP:journals/access/WangBH19}, novelty detection (ND) \citep{DBLP:journals/sigpro/PimentelCCT14}, and out-of-distribution (OOD) detection \citep{DBLP:conf/cvpr/GrahamPTNOC23}.
\citet{DBLP:journals/corr/abs-2110-11334} review the similarities and the subtle differences of these topics.
Although AD, OD, ND, and OOD detection are different topics, these terms (especially AD, OD, and ND) are used interchangeably in some previous works \citep{DBLP:journals/pieee/RuffKVMSKDM21,DBLP:journals/csur/PangSCH21}.
It is worth mentioning that reconstruction-based methods can also be used in ND \citep{DBLP:conf/nips/PidhorskyiAD18,DBLP:conf/cvpr/PereraNX19} and OOD detection \citep{DBLP:conf/cvpr/Zhou22}.
% Strictly speaking, AD, OD, ND, and OOD detection are different topics.
% Nevertheless, these terms, especially AD, OD, and ND, are used interchangeably in some previous work \citep{DBLP:journals/pieee/RuffKVMSKDM21}.
% Moreover,
The analysis in this paper can be easily extended to the reconstruction-based methods for ND and OOD detection.
% In a broad sense, anomalies refer to fraudulent or malicious (instead of novel and interesting in ND) observations that appear to be erroneous (instead of rare in OD) compared to the pre-defined normal data.
% Different from OOD detection, the normal data is often treated as a whole in anomaly detection;
% % no classification inside the normal data is required.

% Previous works on RBAD are mostly empirical studies that focus on boosting the detection efficiency.
% As mentioned in \cref{sec-introduction}, traditional studies on AD have a solid theoretical foundation \citep{DBLP:journals/csur/ChandolaBK09}.
% However, RBAD is among the deep anomaly detection methods that traditional theory cannot explain.
% Despite the empirical success of RBAD, the theory behind this method remains unclear.
% To the best of our knowledge, no previous works have theoretically analyzed the reconstruction-based methods for AD, ND, or OOD detection.

% This paper studies one of the empirical method of  from a theoretical point of view.


% Note that our data model can also be used in the analysis of OOD detection, OD, and ND.
% However, our results are only valid for RBADs since the proceeding feature learning analysis is confined to RBAD.